\section{The Package Documentation}\label{sec:PackageDocumentation}
A \lstinline|package|\index{package} in \Java serves more or less the same purpose as a namespace\index{namespace} in other languages, like Python\index{Python} or \mbox{C++}\index{C++}, but it is not exactly the same. Refer to \autocite{ORACLE_DOC_LANGUAGE_SPECIFICATION:PackageDeclarations} for more on this topic.

Different from the module\index{module}, there is no special source file that defines a \lstinline|package|\index{package}. Instead all source files that are located in the same folder on the directory tree for the source code are in the same \lstinline|package|; this folder structure is maintained for the \texttt{*.class} files, too.

As there is a need for package-level documentation, and because it should be possible to set annotations\index{annotation} to packages\index{package}, the \texttt{package-info.java}\index{package-info.java} file in the package folder is used

\keeplisting{20}
A \texttt{package-info.java}\index{package-info.java} file looks like this:
\begin{lstlisting}[caption={package-info.java}]
/*
 * ==================================================================
 * Copyright © <Year> <Copyright Notice>
 * ==================================================================
 *
 * <License Notice>
 */
 
/**
 *  <package description>
 */
@<AnAnnotation> 
package <packagename>;

import <AnAnnotation>

/*
 *  End of File
 */
\end{lstlisting}

Originally a file named \texttt{package.html}\index{package.html} was used; it contained an HTML\index{HTML} document with the package documentation. Although this still works, you should use not use it anymore, because it does not allow to set annotations on package level. 

The \texttt{<package description>} is a standard Javadoc comment\index{javadoc} comment (refer to the chapter \tqfullvref{sec:DocumentationComments}) and can be as lengthy as necessary. It will be discussed in detail in chapter \ref{sec:PackageComment}.

If the package is annotated\index{annotation}\footnote{I recommend to used the \lstinline|@API| annotation even on the package level (refer to chapter \tqfullvref{sec:APIAnnotation} for details).}, the annotations have to be placed directly before the \bdq{\lstinline|package …|} line (line~13 in the sample above), and the imports for the annotations follow below.

The chapter \tqfullvref{sec:Packages} discusses how to define the name of the package (\texttt{<packagename>}).


\subsubsection{Principles}
\begin{enumerate}[label=P\arabic*.]
	\item{A \texttt{package-info.java}\index{package-info.java} has a beginning and a closing comment like a regular Java source file. See \tqfullref{sec:BeginningComment} and \tqfullref{sec:ClosingComment}.}
	\item{If the \texttt{package-info.java}\index{package-info.java} imports, the \lstinline|import| statements follow the \lstinline|package| statement\index{package}, separated with an empty line.}
	\item{The package description provides complete information about the purpose of the package and its contents.}
\end{enumerate}

%==============================================================================
% End of file!!
%==============================================================================
